% ================================================================
%  conteudo/material_metodos.tex
% ================================================================
\section{Material e Métodos}

\subsection{Material}

\begin{table}[H]
\caption{Materiais utilizados no experimento}
\centering
\begin{tabular}{ll}
\toprule
\textbf{Material} & \textbf{Especificação} \\
\midrule
Uvas                  & Variedade roxa \\
Água peptonada        & Concentração de 0,1\% \\
Meio de cultura       & Sólido \\
Placas de Petri       & 6 unidades \\
Swabs                 & Estéreis \\
Béquer                & -- \\
Bastão de vidro       & -- \\
Bico de Bunsen        & -- \\
Estufa bacteriológica & -- \\
\bottomrule
\end{tabular}
\fonte{Os autores (2026)}
\end{table}

\subsection{Métodos}

O experimento foi conduzido utilizando amostras de uvas (variedade roxa), sob rigorosas condições de assepsia para evitar contaminações externas. Inicialmente, procedeu-se à higienização das mãos e da bancada de trabalho com álcool 70\%, mantendo-se o campo estéril durante todo o processo com o auxílio da chama de um bico de Bunsen. Para a coleta dos microrganismos, foram empregados swabs estéreis, previamente umedecidos em solução de água peptonada a 0,1\%, garantindo a viabilidade celular durante a transferência.

A amostragem foi dividida em duas frentes: a coleta da microbiota presente na superfície externa dos frutos e o isolamento a partir da polpa macerada. No total, foram inoculadas seis placas de Petri, utilizando-se a técnica de esgotamento por estrias, na qual o swab foi deslizado sobre o meio de cultura de forma contínua para promover a separação física das colônias. Após a etapa de inoculação, as placas devidamente identificadas foram transferidas para uma estufa bacteriológica, onde permaneceram incubadas a 30~°C por um período de seis dias, permitindo o desenvolvimento e a posterior análise das culturas isoladas.